The record is deliberately narrow, and the following limitations bound
every claim in this paper.

\paragraph{Scale and stimulus regime.}
All CLLA-era results use one organism size (52 internal/output
neurons), three seeds, 8-channel synthetic Poisson patterns at 20 Hz,
and one operating cell (\(\beta = 0.0046875\), \(\tau_s = 5000\) ms).
Scale studies (E7/E8) exist but were not part of the allocation and
recruitment sequence; no claim about larger substrates is made.

\paragraph{No closed loop.}
Nothing in Phase I closes a behavior loop. E19--E23 showed output
activity is stimulus-locked and (in E23) that consequence shaping
exists but inverts; closed-loop memory behavior is listed as not
established.

\paragraph{Protocol defects bound several verdicts.}
The CLLA capability F2/F6 verdicts are formally unmeasured (empty
frozen windows; adjacent-window diagnostics only). The V2.1 Stage-C
claim is retracted. The fe-era afferent-derived source was vacuous and
was corrected; the recruitment design review's drive\(\to\)posts map
has an unmeasured middle band, so bootstrap magnitudes were bracketed,
not calibrated. Each is reported where it matters.

\paragraph{Operating-point confounds.}
First-exposure burst overshoots (baseline D arm peaks 138.8 Hz;
gain registration pres-1 378--449 Hz) sit near the runaway threshold.
The gain experiments were closed with the operating point not fully
clean under the strict 4b reading, which bounds the strength of the
falsification claim recorded for registration B.

\paragraph{Theoretical claims are minimal.}
The paper does not claim biological fidelity, optimality, or any
statement about what real nervous systems must do; the substrate is
four LIF populations with textbook plasticity, studied as an
experimental system. Interpretive claims (``the shared additive budget
is the interference medium'') are labeled supported-versus-hypothesis
in the text.

\paragraph{What is not claimed about failure.}
``Not supported'' and ``not established'' are never converted into
``impossible.'' In particular: the recruitment-gain hypothesis is
bounded by two registrations at one magnitude; other operating points,
other gain forms, and substrate-preserving mechanisms are untested.
Alternating coexistence being demonstrated does not establish general
memory capability; it establishes coexistence under alternation.